\documentclass[../main.tex]{subfiles}
\begin{document}

\begin{center}
  \Large{\textbf{TÓM TẮT NỘI DUNG}}\\
\end{center}
\vspace{1cm}
Các máy chủ Linux thường trở thành mục tiêu của các hành vi sau xâm
nhập như leo thang đặc
quyền, thiết lập kết nối ngược hay đánh cắp dữ liệu. Vì các hành vi
này đều phải tương tác
với nhân hệ điều hành qua lời gọi hệ thống (syscall), chuỗi syscall
là nguồn quan sát quan
trọng cho các hệ thống phát hiện xâm nhập trên máy chủ (HIDS). Tuy
nhiên, các hướng tiếp
cận hiện tại đối mặt với một khoảng cách thực thi chính sách
(enforcement gap): các mô hình
học sâu như Transformer có khả năng học mẫu tuần tự tốt nhưng quá
nặng để chạy trực tiếp
trong kernel, trong khi các công cụ luật tĩnh như Falco hay Tetragon
có chi phí thấp nhưng
phụ thuộc vào tri thức chuyên gia và khó bao phủ đầy đủ các biến thể tấn công.

Luận văn đề xuất Guepard Shield, một hệ thống HIDS lai lấp khoảng
trống đó bằng cách tách
giai đoạn học biểu diễn ngoại tuyến ra khỏi giai đoạn thực thi thời
gian thực. Hướng tiếp
cận là chưng cất tri thức từ Transformer Teacher sang một automata
hữu hạn đơn định (DFA)
Student đủ nhỏ để nạp vào BPF map và truy vấn với độ phức tạp $O(1)$
mỗi syscall. DFA vừa
phù hợp với ràng buộc của eBPF Verifier, vừa bảo toàn tri thức ngữ
cảnh mà Teacher đã học.

Pipeline gồm bốn giai đoạn: token hóa tên syscall và tạo cửa sổ trượt
(Giai đoạn~1),
huấn luyện Decoder-only Transformer trên dữ liệu bình thường (Giai
đoạn~2), rời rạc hóa
vector trạng thái ẩn bằng K-Means để xây dựng bảng chuyển DFA (Giai
đoạn~3), và nạp bảng
chuyển vào BPF map để thực thi qua eBPF (Giai đoạn~4).

Thực nghiệm trên LID-DS-2021 cho thấy Teacher đạt AUROC = 0.8503 và
PR-AUC = 0.7093 trên
gần 100 triệu window kiểm thử không dùng nhãn tấn công khi huấn
luyện. DFA Student tại
$K = 64$ tạo ra 64 trạng thái và 1.973 bước chuyển (~204~KB), đạt tỷ
lệ phát hiện cấp độ
bản ghi 90.85\%, tỷ lệ báo động giả 1.41\% và fidelity 98.12\% so với
Teacher, xác nhận
tính khả thi của hướng chưng cất Transformer sang DFA cho bài toán
giám sát syscall
thời gian thực.
\begin{flushright}
  \begin{tabular}{@{}c@{}}
    Học viên\\
    \\
    \\
    \\
    \textbf{Nguyễn Minh Đức}
  \end{tabular}
\end{flushright}

\end{document}
